% Nouveau EGA V (cote 25): archive cote 25, site page 125 (PDF page 126).
% typescript. Draft compared: none.
% AI-transcribed by Claude Opus 5.5 (claude-opus-5-5, adaptive thinking, effort high), Message Batches API, batch ega25b (msgbatch_01A7n5YuXubUPNK1zsepi8RG), 2026-11-13; not yet checked by eye.
\scanpage{25/125}
\inserted{compte tenu que l'hyp. sur $f$ implique $f_*(\underline{O}_X)\simeq\underline{O}_S$, d'où $f_*(\underline{O}_X^*)\simeq\underline{O}_S^*$}
\mspage{\hand{28}}

par $\check{P}$ , et associés à $\underline{L}$ (dans le sens de \hand{\S}.1.) correspondent biunivo-%
quement aux\hand{*} homomorphismes \uline{surjectifs} $Q\longrightarrow\check{E}$\hand{,} modulo \inserted{multiplication par des} sections de $\underline{O}_S^{\hand{*}}$ .\hand{\unsure{$\blacktriangledown$}}
\inserted{i.e. \deleted{aux sections du faisceau $\mathcal{H}om_S(Q,\check{E})^*/\underline{O}_S^*$, où le \unsure{premier} * \unsure{indique}}}
\margin{\unsure{Pour faisceaux d'ensembles} \illegible{2} \deleted{\illegible{2}} \unsure{préciser}}
L'existence de $Q$ se ramène par le procédé brévété au cas S noethé-%
rien, et dans ce cas n'est autre que III 7.7.6. (l'hypothèse sur les
fibres de X étant d'ailleurs inutil\unsure{e|è}). Comme $\underline{E}$ est loc libre de type fini
la donnée d'un hom $\underline{E}\longrightarrow f(\underline{L})$ équivaut à celle d'une section de
$\underline{L}\otimes_{\underline{O}_S}\check{\underline{E}}$ , donc à celle d'un homomorphisme $Q\longrightarrow\check{\underline{E}}$ . Il reste à exprimer
que la condition envisagée dans \hand{\S}.1. est bien vérifiée, ce qui (grâce à
l'hypothèse faite sur les fibres de X/S) se réduit à vérifier que
fibre par fibre l'homomorphisme correspondant $E(s)\longrightarrow$ \xout{Q(s)}\hand{$\vee$}
$H^0(X_s,L_s)\simeq\mathrm{Hom}_{k(s)}(Q(s),k(s))$ est injectif, ou encore $Q(s)\longrightarrow\check{E}(s)$ surjec%
tif, ce qui par Nakayama signffie aussi que $Q$ \quad $E$ est surjectif. Le ``mo
dulo sections de $\underline{O}_S$ '' provient du ``modulo isomorphisme'' dans \hand{\S}.1.\hand{/}

On peut \xout{xxxxxxxxx}\inserted[typed]{interpréter} \hand{\S}.1. par une autre voie, en utilisant le fait
que $P(\underline{Q})$ représente le \xout{xxxxxxxx} sous-foncteur de $\mathrm{Div}_{X/S}$ défini par
$\underline{L}$, en v ertu du n$^{\circ}$2. Par suite, un\xout{x}système linéaire de diviseurs paramétr
par $\check{P}$ \inserted[typed]{et associé à $\underline{L}$} s'interprète comme un morphisme $\check{P} = P(\check{\underline{E}})\longrightarrow P(Q)$; le caractère
\uline{linéaire} de la famille de diviseurs définie par un tel morphisme
s'interprète alors par le fait que ce morphisme soit ``linéaire'', i.e.
précisément défini par un morphisme surjectif $Q\longrightarrow\check{\underline{E}}$ . \hand{$\ulcorner$} On voit aussi
dans\xout{xx} ce cas que le morphisme $\check{P}\longrightarrow\mathrm{Div}_{X/S}$ est un monomorphisme \deleted{\illegible{2}}
(puisque $\check{P}\longrightarrow P(Q)\hookrightarrow\mathrm{Div}_{X/S}$ l'est), fait d'ailleurs \inserted{plus} général, cf corol-%
laire plus bas. Convenons \inserted[typed]{alors} de dire que d\unsure{e|a}ux familles linéaires de divi-%
seurs de X/S, paramétrées par les fibrés projectifs $\check{P},\check{P}'$ sont \uline{isomorph-}%
\uline{es} s'ils sont transformés l'un en l'autre par un S-isomorphisme $\check{P}\longrightarrow\check{P}'$ ,
(lequel sera d'ailleurs unique grâce au fait qu'on a des monomorphismes
dans $\mathrm{Div}_{X/S}$ ). \deleted{On appelera système linéaire de diviseurs dans X/S}
(\deleted{sans référence à un fibré projectif précisé qui le paramétriserait) une}

\mspage{\hand{125}}
